\chapter{Cálculo $\lambda$}

\begin{defn}
Utilizamos una representación \textit{locally nameless} 
para el cálculo $\lambda$, donde las variables ligadas se representan 
como números naturales (índices de De Bruijn) y las variables libres como
términos del tipo \lean{Var}. El tipo \lean{Var} representa nombres: 
la igualdad es decidible en ellos (\lean{DecidableEq}) 
y es posible generar uno nuevo para cualquier conjunto finito dado
de nombres (\lean{HasFresh}).
\begin{leancode}
universe u
variable {Var : Type u} [HasFresh Var] [DecidableEq Var]

inductive Term (Var : Type u)
  | bvar : ℕ → Term Var
  | fvar : Var → Term Var
  | abs  : Term Var → Term Var
  | app  : Term Var → Term Var → Term Var
\end{leancode}
Por ejemplo, podemos representar la expresión ``$\lambda x. \, x \, y$'' de la siguiente manera.
\begin{leancode}
variable {y : Var}
def ejemplo_rep := abs (app (bvar 0) (fvar y))
\end{leancode}
Lo cual también se puede escribir como ``$\lambda . \, 0 \, y$''. 
\end{defn}

\begin{defn}
La substitución consiste en remplazar una variable libre por un término. 
\begin{leancode}
def subst (m : Term Var) (x : Var) (sub : Term Var) : Term Var :=
  match m with
  | bvar i  => bvar i
  | fvar x' => if x = x' then sub else fvar x'
  | app l r => app (l.subst x sub) (r.subst x sub)
  | abs M   => abs <| M.subst x sub
\end{leancode}
Definimos una notación para la substitución de variables libres que imita la notación matemática estándar.
\begin{leancode}
scoped notation t:max "[" x ":=" t' "]" => subst t x t'
\end{leancode}
\end{defn}
